\documentclass[11pt,a4paper]{article}

%% PACHETE MATEMATICE
\usepackage{amsmath,amssymb,amsfonts}
\usepackage{amsthm}
\usepackage{mathpazo}
\usepackage{eulervm}
\usepackage{enumerate}
\usepackage{color}
\usepackage{graphicx}
\usepackage[all]{xy}
\usepackage{xcolor}
\usepackage{framed}
\usepackage[unicode]{hyperref}
\setcounter{MaxMatrixCols}{20}
\usepackage{multicol}
%\usepackage[romanian]{babel}


%% ASPECT
\linespread{1.05}
\usepackage[T1]{fontenc}
\usepackage[utf8x]{inputenc}
\usepackage[margin=2cm]{geometry}
% \usepackage[color]{showkeys}
\usepackage{anyfontsize}
\usepackage{titlesec}
\titleformat*{\section}{\large\bfseries}

\usepackage{fancyhdr}
\pagestyle{fancy}
\rhead{\small \color{gray} Notițe de Adrian Manea}
\lhead{\small \color{gray} Aero-BCD, 2017---2018, M. Olteanu \& L. Costache}


\usepackage[autostyle = false, style = german]{csquotes}
\usepackage[romanian]{babel}
\newcommand\qq{\enquote}

%% COMENZILE MELE
\renewcommand*{\proofname}{Dem.:}
\newcommand\bb{\mathbb}
\newcommand\dr{\mathrm}
\newcommand\kal{\mathcal}
\newcommand\fk{\mathfrak}
\newcommand\op{\oplus}
\newcommand\ot{\otimes}
\newcommand\ds{\displaystyle}
\newcommand\id{\indent}
\newcommand\nid{\noindent}
\newcommand\seq{\subseteq}
\newcommand\sneq{\subsetneq}
\newcommand\speq{\supseteq}
\newcommand\spneq{\supsetneq}
\newcommand\rar{\rightarrow}
\newcommand\Rar{\Rightarrow}
\newcommand\xrar{\xrightarrow}
\newcommand\lar{\leftarrow}
\newcommand\Lar{\Leftarrow}
\newcommand\xlar{\xleftarrow}
\newcommand\lrar{\leftrightarrow}
\newcommand\Lrar{\Leftrightarrow}
\newcommand\ideal{\trianglelefteq}
\newcommand\ai{\text{ a.\^{i}. }}
\newcommand\sk{(\textit{Schi\c{t}\u{a}}) }
\newcommand\rhup{\rightharpoonup}
\newcommand\rhdn{\rightharpoondown}
\newcommand\lhup{\leftharpoonup}
\newcommand\lhdn{\leftharpoondown}
\newcommand\vid{\emptyset}
\newcommand\wh{\widehat}
\newcommand\ol{\overline}
\newcommand\NN{\mathbb{N}}
\newcommand\RR{\mathbb{R}}
\newcommand\ZZ{\mathbb{Z}}
\newcommand\QQ{\mathbb{Q}}
\newcommand\CC{\mathbb{C}}

%% STILURILE MELE
\newtheoremstyle{dotless}{}{}{\itshape}{}{\bfseries}{:}{ }{}

\theoremstyle{dotless}
\newtheorem{theorem}{Teorem\u{a}}[section]
\newtheorem{proposition}{Propozi\c{t}ie}[section]
\newtheorem{lemma}{Lem\u{a}}[section]
\newtheorem{corollary}{Corolar}[section]
\newtheorem{exercise}{Exerci\c{t}iu}

\newtheoremstyle{dotlessdr}{}{}{}{}{\bfseries}{:}{ }{}
\theoremstyle{dotlessdr}
\newtheorem{example}{Exemplu}[section]
\newtheorem{definition}{Defini\c{t}ie}[section]
\newtheorem{remark}{Observa\c{t}ie}[section]




\begin{document}

\begin{center}
  {\large \textbf{Seminar 4 \\ Convergență absolută. Calculul sumelor seriilor}}
\end{center}

\vspace{1cm}

1. Studiați convergența absolută și semiconvergența seriilor cu termenul general $x_n$, dat de:
\begin{enumerate}[(a)]
\item $x_n = (-1)^{n+1} \frac{2^n \sin^{2n} x}{n + 1}$;
  \hfill (rădăcină $\Rightarrow 2 \sin^2 x$, discuție, Leibniz);
\item $x_n = \frac{(-1)^{n+1}}{n^{\alpha + \frac{1}{n}}}, \alpha \in \RR$;
  \hfill (comparație 3 $\frac{1}{n^\alpha}$, discuție, Abel);
\item $x_n = \frac{\cos n\alpha}{n^2}, \alpha \in \RR$; \hfill (AC, comparație $\frac{1}{n^2}$);
\item $x_n = (-1)^{n+1} \frac{1}{\sqrt{n(n+1)}}$; \hfill (SC).
\end{enumerate}

\vspace{1cm}

2. Aproximați cu o eroare mai mică decît $\varepsilon$ sumele seriilor definite de
șirul $x_n$:
\begin{enumerate}[(a)]
\item $x_n = \frac{(-1)^n}{n!}, \varepsilon = 10^{-3}$; \hfill (R: $n = 7$);
\item $x_n = \frac{(-1)^n}{n^3 \sqrt{n}}, \varepsilon = 10^{-2}$; \hfill (R: $n = 4$);
\item $x_n = \frac{1}{n!}, \varepsilon = 10^{-3}$; \hfill (R: $n = 6$);
\item $x_n = \frac{1}{n! 2^n}, \varepsilon = 10^{-4}$; \hfill (R: $n = 5$);
\end{enumerate}

\vspace{1cm}

3. Să se demonstreze convergența și să se determine sumele seriilor cu termenul general dat de:
\begin{enumerate}[(a)]
\item $x_n = \frac{1}{(\alpha + n)(\alpha + n + 1)}$; \hfill (R: $\frac{1}{1 + \alpha}$);
\item $x_n = \frac{1}{16n^2 - 8n - 3}$; \hfill (R: $\frac{1}{4}$);
\item $x_n = \ln \frac{n+1}{n}$; \hfill (R: $\infty$).
\end{enumerate}

\vspace{1cm}

4. Să se stabilească natura seriilor cu termenul general dat de:
\begin{enumerate}[(a)]
\item $x_n = \arcsin \frac{n+1}{2n + 3}$; \hfill (D, necesar);
\item $x_n = \frac{\ln n}{2n^3 - 1}$; \hfill (comparație $\frac{n}{n^3}$);
\item $x_n = \frac{1}{n \sqrt[n]{n^3}}$; \hfill (comparație 3 $\frac{1}{n}$);
\item $x_n = e^{-n^2}$;
\item $x_n = ne^{-n^2 + n}$;
\item $x_n = \frac{(an)^n}{n!}, a \in \RR$; \hfill (C $a < \frac{1}{e}$);
\item $x_n = \frac{a^n}{n!}, a > 0$; \hfill (Raabe);
\item $x_n = \Big( \frac{(2n - 1)!!}{(2n)!!} \Big)^2 \cdot \frac{1}{n}$; \hfill (Raabe);
\item $x_n = \Big( \frac{1 \cdot 4 \cdot 7 \cdots (3n - 2)}{3 \cdot 6 \cdot 9 \cdots (3n)} \Big)^2$; \hfill (Raabe);
\item $x_n = \frac{1}{(\ln \ln n)^2}$; \hfill (C, logaritmic);
\item $x_n = \frac{\ln n}{\sqrt[3]{n^4}}$; \hfill (C, logaritmic);
\item $x_n = (-1)^{n+1} \frac{2n+1}{3^n}$; \hfill (C, Leibniz);
\item $x_n = \frac{\cos n \cdot \cos \frac{1}{n}}{n}$; \hfill (C, Abel).
\end{enumerate}




\end{document}
